% =====================================================================
%  Supplementary Notes (SN6): per-paper case reports for the two
%  AI-drafted, curator-reviewed annotations (Duerr 2006, Inouye 2018).
%  Referenced from the main text (sec:examples) as Supplementary Note SN6.
%
%  This note gives the item-level detail the main paper omits for space:
%  a decomposition table, one worked annotation item, and the candidate
%  extensions each paper surfaces. Review verdicts are in SN4 and the
%  term-level crosswalk in SN5; the complete YAML annotations are the
%  repository files annotations/v0/duerr2006.yaml and
%  annotations/v1/inouye2018.yaml. The YAML excerpts below are abridged
%  for readability.
% =====================================================================
\section{Case reports: Duerr 2006 and Inouye 2018}
\label{sec:supp-crep}

The two AI-drafted annotations illustrate contrasting parts of the model's
operating range. The main paper (Section~4) summarizes them; this
note provides the item-level detail omitted there for space. Duerr is the
more conventional case: six of its seven \GeneticEvidence{} items fit the
current schema without strain, and the remaining gap concerns translational
evidence. Inouye tests the limits of the variant-level model: all four items
require forced-fit assignments, which motivate six candidate extensions.
Assertion-by-assertion review verdicts are given in
\cref{sec:supp-ai-review}; the term-level standards crosswalk is in
\cref{sec:supp-crosswalk}; and the complete annotations are the
repository files \fpath{annotations/v0/duerr2006.yaml} and
\fpath{annotations/v1/inouye2018.yaml}.

\subsection{Case Report CR-DU: Duerr et al.\ 2006 (genome-wide association)}
\label{sec:supp-ex-duerr}

Duerr and colleagues identified \textit{IL23R} as an inflammatory bowel
disease gene through a genome-wide association study of non-Jewish patients
with ileal Crohn's disease, followed by replication in a Jewish cohort and
family-based transmission testing in 883 nuclear families. Our annotation
decomposes the paper into seven \GeneticEvidence{} items
(\cref{tab:cr-duerr}).

\begin{table}[H]
  \centering
  \footnotesize
  \caption[Duerr 2006: evidence-item decomposition]{Duerr 2006: the seven
  \GeneticEvidence{} items, with target type, leaf-level Method, and
  Credibility.}
  \label{tab:cr-duerr}
  \begin{tabularx}{\linewidth}{@{}l >{\raggedright\arraybackslash}X >{\raggedright\arraybackslash}p{1.35cm} >{\raggedright\arraybackslash}p{2.5cm} l@{}}
    \toprule
    ID & Claim & Target & Method & Cred. \\
    \midrule
    GE-1 & rs11209026 protective in non-Jewish ileal CD (discovery) & \dimval{Variant} & \dimval{GWAS} & High \\
    GE-2 & rs11209026 protective in Jewish ileal CD (replication) & \dimval{Variant} & \dimval{Association Study} & High \\
    GE-3 & Gln undertransmitted to affected offspring (FBAT) & \dimval{Variant} & \dimval{Transmission Disequilibrium Test} & High \\
    GE-4 & Multiple independent IL23R signals (conditional) & \dimval{Variant} & \dimval{Fine Mapping} & High \\
    GE-5 & No association at IL12RB1/IL23A/IL12B & \dimval{Gene} & \dimval{GWAS} & High \\
    GE-6 & Cited mouse-model and functional support & \dimval{Gene} & \dimval{In Vivo}, \dimval{In Vitro} & High \\
    GE-new-1 & Cited anti-p40 (IL12B) antibody trial & \dimval{Gene} & \dimval{Clinical Evidence} & High \\
    \bottomrule
  \end{tabularx}
\end{table}

Six of the seven items fit the current schema without strain. The seventh,
GE-new-1, is a cited anti-p40 antibody trial added during the
missing-evidence check. It concerns drug-response evidence for which the
model has no fitting genetic Knowledge Domain, revealing a
translational-scope gap. The discovery item, GE-1, is shown in abridged form
below. It illustrates how dimensions, a credibility comment, and
source-anchored assertions are recorded.

\begin{lstlisting}[language=yaml]
id: GE-1
label: "rs11209026 (Arg381Gln) in IL23R confers protection
         against ileal CD in non-Jewish discovery cohort"
knowledge_domain: [POPULATION_GENETICS]
method: [GWAS]
target_type: VARIANT
target: "rs11209026 (c.1142G>A, p.Arg381Gln) in IL23R"
resolution: VARIANT
variant_ascertainment: [OBSERVED_IN_CASES, OBSERVED_IN_CONTROLS]
phenotype_scale: CLINICAL
credibility: HIGH
credibility_comment: >
  Genome-wide significance after Bonferroni correction
  (corrected P = 1.56e-3); large effect (OR = 0.26);
  independent replication follows in GE-2.
candidate_extension_exercised: CE-DU1   # direction of effect
assertions:
  - id: GE-1.A1
    statement: "rs11209026 is genome-wide significantly
                associated with ileal Crohn's disease"
    source_span: {page: 1,
      key_phrase: "rs11209026 (P = 5.05e-9, corrected P = 1.56e-3)"}
  - id: GE-1.A2
    statement: "The Gln allele is protective (OR = 0.26)"
    source_span: {page: 2,
      key_phrase: "glutamine allele appears to protect against
                   development of CD"}
\end{lstlisting}

\paragraph{Candidate extensions.} Curator review identified three candidate
extensions, all provisional under the two-papers promotion rule:

\begin{itemize}
  \item \textbf{CE-DU1, direction of effect:} a first-class dimension
  distinguishing protective from risk-conferring alleles. In this case, the
  protective direction is implicit in an odds ratio of 0.26. The distinction
  may matter for downstream ACMG/AMP-style use.
  \item \textbf{CE-DU3, cross-item synthesis:} a construct for
    higher-order claims that jointly reference several evidence items. Here,
    the relevant claim is the paper's recommendation that blocking the IL-23
    pathway is a rational therapeutic strategy. The issue arises twice, in
    that recommendation and in GE-new-1, and therefore also raises a question
    about the model's translational scope.
  \item \textbf{CE-DU4, conditional-activation scope:} widen the
    activation condition for
    \{\texttt{mode\_of\_\allowbreak inheritance},
    \texttt{mendelian\_\allowbreak segregation}, \texttt{exact\_variant},
    \texttt{subdomain}\} beyond
    \dimval{Human Genetics}~$+$~\dimval{Gene} to also cover
    \dimval{Human Genetics}~$+$~\dimval{Variant}. GE-3, a family-based
    transmission test of a specific variant, illustrates this need. The
    extension would also require a \dimval{Complex} mode-of-inheritance
    value and a family-based subdomain value.
\end{itemize}
A fourth candidate, CE-DU2, proposed a separate dimension for negated
assertions. It was \emph{retracted} on review because an assertion is already
a predicate and can express either absence or presence. This retraction is
the corpus's clearest example of the conservative promotion rule preventing
a premature addition.

\subsection{Case Report CR-IN: Inouye et al.\ 2018 (polygenic risk score)}
\label{sec:supp-ex-inouye}

Inouye and colleagues developed and externally validated a meta-genomic risk
score (metaGRS) for coronary artery disease from 1{,}745{,}180 variants in
482{,}629 UK Biobank participants. The evidential target is a derived,
whole-genome score, and the primary measurements are epidemiological, such
as hazard ratios and C-indices. The paper therefore exposes limitations of
a model whose target types are variant-level.

The annotation reported here is the \texttt{v1} re-annotation produced
under the current protocol. It was developed after the Duerr review; the
original five-item \texttt{v0} draft remains in the repository for
comparison. The \texttt{v1} annotation decomposes the paper into four
claim-coherent items (\cref{tab:cr-inouye}). Each can be represented only
through deliberate placeholder assignments, recorded as forced fits and
candidate extensions rather than left implicit.

\begin{table}[H]
  \centering
  \footnotesize
  \caption[Inouye 2018: evidence-item decomposition]{Inouye 2018
  (\texttt{v1}): four \GeneticEvidence{} items. Each uses the same
  placeholder Target Type and Resolution value, \dimval{Variant}, and is
  flagged against CE-IN1 and CE-IN2. The actual target is a whole-genome
  polygenic score.}
  \label{tab:cr-inouye}
  \begin{tabularx}{\linewidth}{@{}l >{\raggedright\arraybackslash}X p{1.5cm} l@{}}
    \toprule
    ID & Claim & Target$^{\dagger}$ & Cred. \\
    \midrule
    GE-1     & Prediction and quintile stratification (HR 1.71/SD; 4.17-fold) & \dimval{Variant} & High \\
    GE-2     & Independence from six conventional risk factors & \dimval{Variant} & High \\
    GE-3     & Stratification persists among medicated individuals & \dimval{Variant} & High \\
    GE-new-1 & Combining three component scores outperforms each & \dimval{Variant} & Medium \\
    \bottomrule
  \end{tabularx}

  {\footnotesize $^{\dagger}$Placeholder assignment. The actual target is a
  polygenic \dimval{Score} (CE-IN1) at whole-genome-aggregate resolution
  (CE-IN2). Method is recorded at the intermediate
  \dimval{Association Study}/\dimval{Meta-Analysis} level because no leaf
  represents polygenic-score association (flag F-METH). Variant Ascertainment
  is marked not applicable, which supports the CE-IN1 SHACL argument.}
\end{table}

The independence item, GE-2, makes these limitations concrete and is shown
in abridged form below. Protocol v0 (Duerr) used the field name
\texttt{resolution}; protocol v1 (Inouye) renamed it
\texttt{target\_resolution}. The corpus retains both spellings, and the
YAML-to-RDF transformation normalizes both to \texttt{gem:resolution}, so
\texttt{gem-validate} treats them identically. \texttt{target\_type} and
\texttt{target\_resolution} are assigned the placeholder value
\dimval{Variant} and flagged against CE-IN1 and CE-IN2. Variant
Ascertainment is marked not applicable because it does not apply cleanly to
a score target. The epidemiological measures are recorded in
\texttt{special\_considerations} and motivate CE-IN4.

\begin{lstlisting}[language=yaml]
id: GE-2
label: "metaGRS predicts incident CAD largely independently
        of six conventional risk factors"
knowledge_domain: [POPULATION_GENETICS]
method: [ASSOCIATION_STUDY]    # intermediate; no leaf fits PGS assoc. (F-METH)
target_type: VARIANT           # forced; CE-IN1 proposes SCORE
target_resolution: VARIANT     # forced; CE-IN2 -> WHOLE_GENOME_AGGREGATE
variant_ascertainment: not_applicable_or_omitted  # CE-IN1 (SHACL); F-VA
phenotype_scale: CLINICAL
credibility: HIGH
# epidemiological measures (C-index, HR) in special_considerations; CE-IN4
assertions:
  - id: GE-2.A1
    statement: "metaGRS has a higher C-index than any single factor"
    source_span: {page: 6,
      phrase: "higher C-index (C = 0.623; 95% CI: 0.615 to 0.630)"}
  - id: GE-2.A3
    statement: "Adding metaGRS to all six factors raised C-index to 0.696"
    source_span: {page: 6,
      phrase: "C-index of 0.696 (95% CI: 0.688 to 0.703)"}
\end{lstlisting}

\paragraph{Candidate extensions.} The Inouye annotation identified six
candidate extensions, the most of any paper in the corpus:

\begin{itemize}
  \item \textbf{CE-IN1, \dimval{Score} target type:} a metaGRS is
  neither a gene nor a variant, but a derived composite predictor. A
  \dimval{Score} value, possibly characterized as polygenic, composite, or
  both, would represent the target directly. Under the current
  \dimval{Variant} placeholder, all four items mark Variant Ascertainment
  as not applicable. The schema accepts this as an explicit escape, but a
  first-class \dimval{Score} target would remove the forced fit.

  \item \textbf{CE-IN2, \dimval{Whole Genome Aggregate} resolution:} the
  score aggregates variants across the genome non-contiguously, which the
  current contiguous or pointwise Resolution values do not capture.

  \item \textbf{CE-IN3, target-composition dimension:} an axis orthogonal
  to Resolution that distinguishes \dimval{Single},
  \dimval{Aggregate}, and \dimval{Composite} targets. This provides a
  complementary framing of the same underlying gap as CE-IN2.

  \item \textbf{CE-IN4, epidemiological measurement targets:} an
  extension of \texttt{measurement\_target}, which is currently oriented
  toward molecular assays, with values such as \dimval{Hazard Ratio},
  \dimval{Odds Ratio}, \dimval{C-Index}, and \dimval{AUROC}. These values
  would be activated for population-genetics evidence.

  \item \textbf{CE-IN5, structured cohort descriptor:} a replacement for
  free-text cohort properties with named fields for size, ancestry
  composition, follow-up, and ascertainment. Such a structure could be
  reused across Duerr, Davis, and future cohort studies.

  \item \textbf{CE-IN6, natural versus derived-artifact target:} a flag,
  or a \dimval{Natural}/\dimval{Derived} dimension, indicating that the
  metaGRS is a human-designed construct rather than a natural kind.
\end{itemize}

